\section{Following An Envelope}\label{sec:envelope}

Follow how a sound develops rather than only how bright its final spectrum
looks. Use an oscillator or noise source, a VCA, and a repeating envelope.

\paragraph{Starting point}
Load \textbf{DecibelInspection}. It sets unity input gain, Flattop, no
smoothing, zero Slope, and a $-90$ to $+12\,\mathrm{dB}$ color range.
Change Window to Hann. Use Linear Freq Scale and $0$--$5\,\mathrm{kHz}$
for a pitched source, or a wider range for noise. Keep Rack at
$48\,\mathrm{kHz}$ for the timing examples here.

\begin{enumerate}
  \item Patch the sound through the VCA. Send the envelope to the VCA's
  level input, and branch the VCA output to Spectre and your listening path.
  Trigger a short attack and a decay of roughly half a second once a second.
  \item Watch several repetitions. A sine leaves one band whose color
  brightens at each attack and fades during decay. Noise leaves a broad
  vertical event. Your envelope shape controls how those colors fade.
  \item Set Average to $300\,\mathrm{ms}$ and observe new events. Their
  level changes should look more gradual. Previously recorded columns are
  unchanged; Average affects newly analyzed spectra. Restore Average to zero.
  \item Send an envelope to oscillator pitch or to a filter's cutoff.
  Compare movement in frequency with fading amplitude. Freeze after a few
  events and inspect their shapes; resume to try a different envelope.
\end{enumerate}

\paragraph{What time detail survives?}
At $48\,\mathrm{kHz}$ each column summarizes about $42.7\,\mathrm{ms}$
of audio, and successive windows overlap by half. A very short click can
therefore affect several columns. Its drawn width is not its exact duration.
Averaging adds further persistence. Use an oscilloscope for individual sample
edges or precise attack timing, and Spectre for the changing frequency content.

\paragraph{Try a fair comparison}
Keep gain, Floor, Ceil, window, and sample rate fixed while changing the
source. A darker decay after raising Floor does not show that the VCA became
quieter: it shows that you hid more low-level detail. Likewise, changing
Slope can make upper partials appear to sustain longer by making them more
visible. Use zero Slope while comparing decay across frequencies.
